\section*{\S\ 5. Die Zerlegungsidentit\"aten in expliziter Form.}
Spezialisieren wir in (25.) im Falle 4) $\tau$ zu $\sigma+\rho$, so kommt:
\begin{equation}
 \pair{q_\rho A^\sigma}{y^{(1)}\cdots y^{(\rho+\sigma)}}
 =\sum_{i_\rho=i_{\rho-1}+1}^{\tau}
 \eps_{i_\rho}\pair{q_\rho}{y^{(i_1)}\cdots y^{(i_\rho)}}
 \pair{A^\sigma}{y^{(i_{\rho+1})}\cdots y^{(i_{\rho+\sigma})}},
\tag{32}
\end{equation}
wo $\eps_{i_\rho}=(-1)^{\delta_{i_\rho}}$, $\delta_{i_\rho}=\rho(\rho+1)/2+i_1+\cdots+i_\rho$ wird. Dies ist eine allgemeinere Form des Produktsatzes f\"ur Matrizen als (16.) und (17.) und ergibt sich durch Laplacesche Entwicklung der Determinante der Produktenmatrix
\[
 \sum\pm(v^{(1)}y^{(1)})\cdots(v^{(\rho)}y^{(\rho)})(a^{(1)}y^{(\rho+1)})\cdots(a^{(\sigma)}y^{(\rho+\sigma)}).
\]
Der allgemeinere Produktsatz ist also aus den Identit\"aten (20.), bzw. (21.) zusammengesetzt, und das Herausgreifen der speziellen Form (16.), (17.) ist dadurch erkl\"art.

Die Relation (32.) gibt nun das Mittel zur expliziten Darstellung der Zerlegungsidentit\"aten. Wir betrachten dazu die in (31.) auftretenden Formen als Grundformen; d. h. wir f\"uhren die Substitution (11.) aus, die an der expliziten Darstellung in den Reihen $A,B$ nichts \"andert, da nach (8.) die neu eingef\"uhrten Reihen $R$ sich durch die Reihen $A,B$ selbst und nicht durch deren Teilreihen $a^{(1)}a^{(2)}\ldots$, $b^{(1)}b^{(2)}\ldots$ ausdr\"ucken lassen. Sei also
\begin{equation}
 \pair{q_{\rho-\alpha}B^{n-\rho-\sigma+\tau}}{A_{n-\sigma}p_{\tau-\alpha}}
 =\bigl(R_{\rho-\alpha}p_{n-\rho+\alpha}\bigr)
  \pair{R^{\tau-\alpha}}{p_{\tau-\alpha}},
\tag{33}
\end{equation}
so kommt f\"ur die dem Index $\alpha$ entsprechende Einzelsumme von (25.):
\begin{equation}
\begin{aligned}
&\sum \eps_{i_\alpha}\bigl(R_{\rho-\alpha}p_{n-\rho}y^{(i_1)}\cdots y^{(i_\alpha)}\bigr)
 \pair{R^{\tau-\alpha}}{y^{(i_{\alpha+1})}\cdots y^{(i_\tau)}}\\
&\quad=\sum \eps_{i_\alpha}\bigl([R_{\rho-\alpha}p_{n-\rho}]^\alpha y^{(i_1)}\cdots y^{(i_\alpha)}\bigr)
 \pair{R^{\tau-\alpha}}{y^{(i_{\alpha+1})}\cdots y^{(i_\tau)}}\\
&\quad=\eps\pair{[R_{\rho-\alpha}p_{n-\rho}]^\alpha R^{\tau-\alpha}}{p_\tau}
 =\eps\pair{R^{\tau-\alpha}q_{n-\tau}}{R_{\rho-\alpha}p_{n-\rho}}.
\end{aligned}
\tag{34}
\end{equation}
Damit ist nun in der Tat ein expliziter Schlu\ss ausdruck gegeben, und die symmetrische Form, in der $q_\rho$ und $p_\tau$ eingehen, zeigt, da\ss\ (30.) zum gleichen Schlu\ss ausdruck f\"uhrt, der also unabh\"angig ist von der urspr\"unglichen Zerlegung. Ein Unterschied zwischen (25.) und (30.) besteht nur, solange die Reihen $y$ und $v$ selbst\"andige Bedeutung haben, die Zerlegungsidentit\"at also nach unserer Definition in eine zerlegbare Identit\"at \"ubergeht. Um zu Identit\"aten mit mehr Symbolreihen zu gelangen, hat man nach \S\ 3 (Schlu\ss) nur eine Reihe $q_\rho$ in $q_{\rho_1}C^{\sigma_1}$, bzw. $p_\tau$ in $p_{\tau_1}C_{\tau-\tau_1}$ aufzul\"osen. Die dann entstehenden Ausdr\"ucke:
\begin{equation}
 \pair{q_{\rho_1}C^{\sigma_1}}{[R^{\tau-\alpha}q_{n-\tau}]_\lambda R_{\rho_1+\sigma_1-\alpha-\lambda}},
 \qquad
 \pair{[R_{\rho-\alpha}p_{n-\rho}]^\lambda R^{\tau-\alpha}}{p_{\tau_1}C_{\tau-\tau_1}},
\tag{35}
\end{equation}
sind nun genau vom Typus der bei zwei Symbolreihen auftretenden, lassen also bei einer (33.) entsprechenden Substitution wieder explizite Darstellung zu, so da\ss\ sich durch Wiederholung explizite Darstellung bei beliebig vielen Reihen ergibt. Bemerkt mag noch werden, da\ss\ je die erste und letzte Summe in (25.) bzw. (30.) die explizite Darstellung ohne Anwendung von (33.) allein durch die Formel (32.) ergibt, oder sich \"uberhaupt auf ein einziges, alle Reihen explizit enthaltendes Glied beschr\"ankt. Die durch Aufl\"osung von $q_\rho$ in $q_{\rho_1}C^{\sigma_1}$ usf. aus (32.) entstehende Relation kann als der Produktsatz in seiner allgemeinsten Form aufgefa\ss t werden.

Zusammenfassend werden wir zeigen:

\emph{Satz III:} Mit den Identit\"aten
\begin{equation}
 \pair{q_\rho A^\sigma}{p_\tau B_{\rho+\sigma-\tau}}
 =\sum_{\substack{0,\tau-\sigma}}^{\substack{\tau,\rho}}
 \eps\pair{R^{\tau-\alpha}q_{n-\tau}}{R_{\rho-\alpha}p_{n-\rho}},
\tag{36}
\end{equation}
wo die Reihen $R$ durch (33.) bestimmt sind, und mit den durch Aufl\"osung von $q_\rho,p_\tau$ in $q_{\rho_1}C^{\sigma_1}$ usf. daraus entstehenden ist die Gesamtheit der unzerlegbaren Identit\"aten ersch\"opft.

In der Tat sind die entstehenden Identit\"aten die allgemeinsten ihrer Art, da die einzelnen Reihen in bezug auf Gewicht und Anzahl keiner Beschr\"ankung mehr unterliegen, wie dies noch bei (20.), (21.) der Fall war. Es existieren aber auch keine weiteren Identit\"aten zwischen diesen Reihen; denn bei Aufl\"osung einer Reihe $p_\tau$ in $y^{(1)}\cdots y^{(\tau)}$ sind die allgemeinsten Identit\"aten aus (20.), also nach \S\ 3 (Schlu\ss) aus (20a.) zusammengesetzt; und zwar ist die im Anschlu\ss\ an (16.) diskutierte Zusammensetzung genau die in \S\ 4 durchgef\"uhrte. Den \"Ubergang zu beliebig vielen Reihen liefert aber, wie die Diskussion von (20a.) zeigt, die Anwendung immer derselben Operation auf die durch Aufl\"osung von $q_\rho$ in $q_{\rho_1}C^{\sigma_1}$ usf. entstehenden Ausdr\"ucke. Es folgt daraus auch, da\ss\ der direkte \"Ubergang von (20.) an Stelle von (20a.) nichts Neues liefert; dies l\"a\ss t sich rechnerisch leicht best\"atigen, indem man einmal eine Reihe $q_\rho$ in (25.) spezialisiert, das andere Mal nach der Methode von \S\ 4 den \"Ubergang von (20.) macht; es entstehen beide Male die gleichen Summen, die durch Anwendung des allgemeinen Produktsatzes (siehe oben) in die aus (36.) sich ergebenden Identit\"aten \"ubergehen. Die Identit\"aten (20.), (21.) entsprechen den Spezialf\"allen $\tau=1$, $\rho=1$; es treten dann nur je zwei Einzelsummen auf, also einer obigen Bemerkung entsprechend explizite Darstellung ohne die Substitution (33.), was mit dem fr\"uheren \"ubereinstimmt.

Zum Schlu\ss\ seien noch zwei Spezialf\"alle der Zerlegungsidentit\"at kurz erw\"ahnt. Setzen wir in (31.): $\tau=\sigma$, $\rho=n-\sigma$ und bezeichnen die Reihe $p_\tau$ als $p'_\sigma\sim q'_{n-\sigma}$, so kommt:
\begin{equation}
\begin{aligned}
\pair{A^\sigma}{p_\sigma}\pair{B^\sigma}{p'_\sigma}
-\pair{A^\sigma}{p'_\sigma}\pair{B^\sigma}{p_\sigma}
=\sum_{\alpha=1}^{\sigma,n-\sigma}
\Pi\pair{q_{n-\sigma-\alpha}B^\sigma}{A_{n-\sigma}p_{\sigma-\alpha}}.
\end{aligned}
\tag{37}
\end{equation}
Das ist die Clebschsche Formel\footnote{Clebsch, a. a. O., S. 25--32. Clebsch weist nach, da\ss\ die einzelnen Glieder der rechten Seiten nur von den Reihen $A,B$ abh\"angen; die explizite Darstellung w\"urde sich durch Anwendung von (32.) auf seine Ausdr\"ucke $\Phi_h$ ergeben.} zur Entwicklung einer Form mit zwei kogredienten Reihen $p_\sigma,p'_\sigma$ nach Elementarkovarianten. Die zu einer Form $\pair{A^\sigma}{p_\sigma}^m\pair{B^\sigma}{p'_\sigma}^n$ geh\"orenden Elementarkovarianten sind dann die folgenden:
\begin{equation}
 \left\{\sum_{\alpha=1}^{\sigma,n-\sigma}
 \pair{q_{n-\sigma-\alpha}B^\sigma}{A_{n-\sigma}p_{\sigma-\alpha}}
 \right\}^{\rho}
 \pair{A^\sigma}{p_\sigma}^{m-\rho}\pair{B^\sigma}{p'_\sigma}^{n-\rho},
 \qquad (\rho=0,1,\ldots,m,n).
\tag{38}
\end{equation}
Sie gehen, wie wir kurz sagen wollen, aus der Grundform durch ,,kogrediente Faltung vom Defekt $\rho$`` (vgl. \S\ 2) hervor.

Setzen wir zweitens: $\tau=\sigma-\lambda$, $\rho=n-\sigma-\lambda$; $B^\sigma=A^\sigma$, so kommt:
\begin{equation}
\begin{aligned}
\pair{q_{n-\sigma-\lambda}A^\sigma}{A_{n-\sigma}p_{\sigma-\lambda}}
&=(-1)^\lambda\pair{q_{n-\sigma-\lambda}A^\sigma}{A_{n-\sigma}p_{\sigma-\lambda}}\\
&\quad+(-1)^{(n-\sigma)(\sigma-\lambda)}
 \sum_{\alpha=1}^{\sigma-\lambda,n-\sigma-\lambda}
 \Pi\pair{q_{n-\sigma-\lambda-\alpha}A^\sigma}{A_{n-\sigma}p_{\sigma-\lambda}}.
\end{aligned}
\tag{39}
\end{equation}
Es sind also, wie in \S\ 2 bemerkt, die die Symbolreihen charakterisierenden Bedingungen (12.) f\"ur ungeraden Defekt $\lambda$ eine Folge der Bedingungen von h\"oherem Defekt. F\"ur eine Linearform $\pair{A^\sigma}{p_\sigma}=\pair{B^\sigma}{p_\sigma}$ ergibt sich aus (39.), da\ss\ sich ihre irreduzibeln invarianten Bildungen, die in den Koeffizienten quadratisch sind, auf solche von geradem Defekt reduzieren; das stimmt mit der bekannten Tatsache \"uberein, da\ss\ eine Linearform im bin\"aren und tern\"aren Gebiet keine invarianten Bildungen besitzt.

Bemerkt sei, da\ss\ die allgemeinen Identit\"aten urspr\"unglich abgeleitet sind unter der Annahme, da\ss\ die einzelnen Reihen den Bedingungen (12.) gen\"ugen. Da indes diese einzelnen Reihen nur linear in die Identit\"aten eingehen, so gelten letztere auch f\"ur die allgemeineren, den Bedingungen (12.) nicht gen\"ugenden Reihen.
